\bookchapter{Chapitre 9 : Travail et apprentissage mobiles}

\chapterintro{Je pensais aussi autrefois que travailler sur smartphone c'était « se débrouiller », jusqu'au jour où j'ai oublié mon ordinateur en déplacement et j'ai réussi à modifier un appel d'offres et signer un contrat uniquement avec mon smartphone. Honnêtement, le goulot d'étranglement du travail mobile n'est souvent pas le smartphone lui-même, mais votre méconnaissance de ses capacités.}

\booksection{Traitement de documents : éditer et créer au creux de la main}

Éditer des documents sur l'écran d'un smartphone, ça ne fait pas la même sensation que sur ordinateur. Le plus gros problème est le petit écran, les doigts épais, la mise en page qui dérape facilement. Mais quelques astuces améliorent nettement l'expérience :

\begin{itemize}
  \item \textbf{WPS Office}: La meilleure localisation, parfaitement compatible avec les formats Microsoft (docx/xlsx/pptx), l'ensemble d'outils PDF intégré est très riche. J'aime particulièrement sa fonction « commentaire vocal » — quand vous relisez un contrat et que taper n'est pas pratique, dictez simplement vos remarques au smartphone, il les convertit automatiquement en annotation texte collée dans le document. De plus, WPS intègre une vaste bibliothèque de modèles gratuits, ce qui fait gagner beaucoup de temps pour les CV et les présentations de rapport.
  \item \textbf{Microsoft 365}: Si votre entreprise utilise Office authentique, c'est le meilleur choix. La compatibilité de format avec la version bureau est quasi parfaite, sans les légers décalages de mise en page que WPS peut parfois produire sur certaines mises en page complexes. Intégration profonde avec OneDrive : le document que vous modifiez à moitié sur l'ordinateur s'ouvre dans sa dernière version sur le smartphone.
  \item \textbf{Shimo Documents/Tencent Documents}: Ce sont les représentants de la « collaboration en ligne ». Adaptés à l'édition simultanée d'un document par plusieurs personnes, les modifications de chacun s'affichent en temps réel. J'ai utilisé Shimo Documents avec des amis pour préparer un guide de voyage, deux personnes remplissant le tableau en même temps, bien plus efficace que de le faire seul.
\end{itemize}

\textbf{Flux de travail pratique}: Mon habitude actuelle — les fichiers importants sur OneDrive, l'édition quotidienne avec WPS, la collaboration d'équipe sur Shimo Documents. Depuis que mes fichiers sont dans le cloud, je n'ai plus jamais eu la gêne du « clé USB oubliée ».

\booksection{Applications de prise de notes : de la note rapide au système de connaissances}

Pour prendre des notes, beaucoup de gens veulent d'abord mettre en place un système complexe, puis abandonnent au bout de trois jours. Mon conseil : commencez léger — adoptez d'abord l'habitude de « noter », puis réfléchissez à la « gestion ».

\begin{itemize}
  \item \textbf{Notes rapides — le système intégré suffit}: L'application Notes iOS se lance extrêmement vite, et prend en charge le scan de documents, l'écriture manuscrite, l'enregistrement vocal. Les applications de notes de chaque marque Android ne sont pas en reste, elles se lancent rapidement même depuis l'écran de verrouillage. Ces outils sont parfaits pour noter des idées temporaires, des listes de courses, des points de réunion. La clé est la rapidité de lancement, vous notez dès que l'idée vient, sans attendre le chargement de l'application.
  \item \textbf{Lecture différée — collecter l'information}: La fonction « Fenêtre flottante » de WeChat est en fait un outil caché. Vous voyez un bon article sans le temps de le lire, cliquez sur « Fenêtre flottante » en haut à droite, il se réduit en un petit point collé au bord de l'écran, accessible à tout moment. Bien plus pratique que de transférer vers « Assistant de transfert de fichiers ». De plus, des outils de lecture différée comme Cubox ou Pocket peuvent collecter sur toutes les plateformes les bons contenus trouvés n'importe où.
  \item \textbf{Construire un système — Notion et Feishu}: Quand vous trouvez que vos notes éparses deviennent difficiles à retrouver, passez à Notion. C'est vraiment flexible — vous pouvez créer une « base de données de notes de lecture », organisée automatiquement par titre, auteur, note ; ou construire un « tableau de projet » visualisant la progression des tâches. La « table multidimensionnelle » de Feishu est similaire, avec un accès plus rapide en Chine. J'utilise Notion pour gérer tous mes supports d'écriture, notes d'apprentissage et progression de projets, depuis plus de trois ans.
\end{itemize}

\textbf{Une stratégie de prise de notes pratique}: Mon habitude personnelle est le « trépied » — idées temporaires dans les notes système (lancement instantané), articles à lire dans la fenêtre flottante WeChat (pratique), connaissances à accumuler sur le long terme dans Notion (systématisé). Ne cherchez pas d'emblée à tout mettre dans un seul outil, c'est le meilleur moyen d'abandonner.

\booksection{Numérisation et OCR : faire entrer le monde papier dans le smartphone}

La numérisation par smartphone, bien maîtrisée, fait vraiment gagner beaucoup de trajets au magasin d'imprimerie. La dernière fois que j'ai fait un visa, j'avais besoin de versions électroniques de plein de documents, tout a été fait au smartphone.

\begin{itemize}
  \item \textbf{Utilisateurs iPhone}: La fonction de scan intégrée à Notes iOS est déjà très bonne. Ouvrez Notes touchez l'icône appareil photo sélectionnez « Scanner un document », visez le document, il détecte automatiquement le cadre et corrige la perspective. Après la prise de vue, vous pouvez ajuster les couleurs — le mode « Noir et blanc » donne un résultat proche d'un scanner, « Niveaux de gris » convient aux documents avec des images.
  \item \textbf{Utilisateurs Android}: Beaucoup d'appareils photo de marque intègrent un mode « Scan de document ». Si vous ne le trouvez pas, téléchargez Microsoft Lens depuis la boutique d'applications (gratuit, et intégré à OneDrive).
  \item \textbf{OCR reconnaissance de texte}: Pour extraire le texte d'un scan pour l'éditer, il faut l'OCR. La fonction « Extraire le texte » intégrée à WeChat est la plus simple — appuyez longuement sur l'image pour l'utiliser, le taux de reconnaissance est correct. Pour un traitement en masse, utilisez l'application Baimiao ou Quanneng Saomiao Wang. Pour la reconnaissance du chinois, le moteur OCR de Baidu donne de bons résultats ; la fonction image vers texte de WPS est également très précise.
  \item \textbf{Astuce de scan}: Placez le document dans un endroit bien éclairé, idéalement avec une lampe de bureau éclairant de côté, pour éviter que votre ombre ne se projette sur le document. Si le document est corné, aplatissez-le avant de prendre la photo. Les documents de plusieurs pages peuvent être scannés en continu puis fusionnés en un seul PDF, pas besoin de les enregistrer page par page.
\end{itemize}

\booksection{Comment choisir un cloud : pas le meilleur, seulement le plus adapté}

J'ai longtemps réfléchi au choix du cloud, et j'ai finalement compris que tout dépend de votre « écosystème principal ».

\begin{itemize}
  \item \textbf{Écosystème Apple complet iCloud}: Synchronisation si fluide qu'on ne la sent pas. Les photos prises sur le smartphone sont automatiquement sur l'ordinateur quelques minutes plus tard. Et Apple fait la protection de la vie privée la plus stricte. Mais les inconvénients sont évidents : le partage interplateforme vers Android ou Windows est compliqué, seulement 5 Go gratuits, l'extension est payante.
  \item \textbf{Windows principal OneDrive}: Si vous utilisez un PC Windows + Office 365, OneDrive est le compagnon idéal. Un document enregistré dans Word est automatiquement dans le cloud. 1 To avec l'abonnement Office est très avantageux. La sauvegarde automatique des photos par l'application smartphone est aussi très pratique.
  \item \textbf{Partage en Chine Baidu Netdisk}: Honnêtement, la limitation de vitesse de téléchargement de Baidu Netdisk est toujours critiquée, mais son écosystème de partage est imbattable. Une grande partie des ressources d'apprentissage et de divertissement sur le web chinois est partagée via Baidu Netdisk. Si vous ne faites que télécharger occasionnellement des fichiers partagés, la version gratuite suffit. Si vous êtes un collectionneur intensif de ressources, prenez l'abonnement, le temps économisé vaut bien plus cher.
  \item \textbf{Alibaba Cloud Drive / Quark Cloud Drive}: Ce sont les nouveaux venus de ces dernières années. Alibaba Cloud Drive n'a pas de limitation de vitesse d'upload/download (pour l'instant), Quark Cloud Drive est profondément intégré au navigateur Quark, vous pouvez enregistrer directement dans le cloud en cherchant des ressources, très pratique.
\end{itemize}

\textbf{Ma solution combinée}: iCloud pour les photos personnelles et les fichiers quotidiens (priorité vie privée), Baidu Netdisk pour les ressources et la réception de partages (priorité écosystème), OneDrive pour les documents professionnels (priorité intégration Office). Trois clouds, chacun son rôle, sans interférence.

\booksection{Un vrai scénario de travail mobile}

Laissez-moi raconter une expérience personnelle. Une fois en déplacement, le client a demandé une modification urgente du contrat. Mon flux de travail :

\begin{enumerate}
  \item À l'hôtel, réception des modifications du client via WeChat (message texte + une photo d'annotations manuscrites).
  \item Ouverture d'OneDrive pour trouver le contrat original, édition avec WPS, modification des clauses en suivant les messages WeChat du client.
  \item Après modification, export en PDF via la fonction « Exporter en PDF » de WPS, envoi dans le groupe WeChat du client pour confirmation.
  \item Le client a proposé une deuxième série de modifications, j'ai utilisé l'outil crayon de WPS pour encercler directement les emplacements sur le PDF, avec des commentaires vocaux pour expliquer la logique des changements.
  \item Après confirmation, signature électronique via la fonction « Signature électronique » de WPS, envoi à l'autre partie pour contre-signature.
  \item Enfin, archivage du contrat signé en PDF dans Baidu Netdisk et iCloud, avec une entrée dans le journal de travail Notion.
\end{enumerate}

Tout le processus a duré environ quarante minutes, sans jamais ouvrir l'ordinateur. Honnêtement, il y a trois ans je ne croyais pas qu'un smartphone pouvait faire tout ça, mais maintenant c'est vraiment possible. La clé est de connaître ces outils à l'avance, ne pas attendre d'en avoir besoin pour télécharger et s'inscrire en panique.

\begin{tcbpractice}{Essayez : construisez votre flux de travail mobile}
\begin{enumerate}
  \item Choisissez votre application document principale (WPS ou Microsoft 365), connectez-vous et associez le cloud.
  \item Ouvrez Notes ou le mémo, scannez un document avec la fonction scan, exportez en PDF et envoyez-le vous.
  \item Choisissez une application cloud, configurez la sauvegarde automatique des photos du smartphone (choisissez « Sauvegarde uniquement en Wi-Fi » pour économiser les données).
  \item Créez une nouvelle page dans Notion ou Feishu, essayez de construire un modèle simple de « journal de travail ».
\end{enumerate}

\end{tcbpractice}
